\documentclass[11pt]{article}
\usepackage[T1]{fontenc}
\usepackage[utf8]{inputenc}
\usepackage{lmodern}
\usepackage[margin=1in]{geometry}
\usepackage{amsmath,amssymb,amsthm,mathtools}
\usepackage{enumitem,booktabs,tabularx,array}
\usepackage{microtype}
\usepackage[hidelinks]{hyperref}
\usepackage{tikz-cd}
\newtheorem{definition}{Definition}[section]
\newtheorem{target}[definition]{Proposed target}
\newtheorem{test}[definition]{Required test}
\theoremstyle{remark}
\newtheorem{remark}[definition]{Remark}
\newcommand{\Obs}{\operatorname{Obs}}
\newcommand{\Real}{\operatorname{Real}}
\newcommand{\Cost}{\operatorname{Cost}}
\newcommand{\im}{\operatorname{im}}
\newcommand{\rank}{\operatorname{rank}}
\newcommand{\conv}{\operatorname{conv}}
\newcommand{\spanlin}{\operatorname{span}}
\newcommand{\RR}{\mathbb R}
\newcommand{\NN}{\mathbb N}
\newcommand{\calE}{\mathcal E}
\newcommand{\calN}{\mathcal N}
\newcommand{\calR}{\mathcal R}
\newcommand{\calB}{\mathcal B}
\newcommand{\calC}{\mathcal C}
\newcommand{\calW}{\mathcal W}
\newcommand{\calU}{\mathcal U}
\newcommand{\calP}{\mathcal P}
\newcommand{\one}{\mathbf 1}
\newcommand{\Id}{\mathrm{id}}
\setlist{nosep,leftmargin=*}
\setlength{\parskip}{3pt}
\title{A research contract for categorical phenomenology\\\large Provisional for discussion}
\author{Revised discussion draft}
\date{3 November 2026}
\begin{document}
\maketitle
\begin{abstract}
The question is how structure emerges from observations, and how different constructions of that structure agree, compare in one direction, lose information, or resist comparison. This provisional contract has two coupled objectives. The comparative programme must explain relations among genuinely different reconstruction methods. The generative programme must construct or explain the selection of structural output that was not supplied in advance. Neither objective can substitute for the other.

Two-sided, context-sensitive presentations remain a strong candidate organizing language. Stacey's comparative method supplies a substantive precedent; the existing generated-image and cylinder results are substantial dependencies. They do not predetermine the answer. Learning dynamics is a prominent candidate generative test: one must explain response-relevant structure created, exposed, or selected during learning, including what the evidence and admissible experiments can distinguish. An unsupplied interface or decomposition remains a demanding flagship, but is not the definition of emergence.

The revision separates available evidence, currently admissible experiments, exact ideal laws for those experiments, and hypothetical laws for a larger repertoire. It retains the earlier positive native-realization and multistage targets, while requiring explicit ambiguity, causal boundaries, nonvacuous input criteria, and falsifying controls. All new theorem statements below are research targets. This document revises a discussion proposal; it is not an adopted contract, a new research schedule, or a claim that these targets have been achieved.
\end{abstract}
\tableofcontents

\section{The question and the proposed direction}
\label{sec:question}

An incidence relation can generate concepts without supplying their names. Complete future laws can generate predictive states and transitions without supplying hidden states. A learning process can reveal that present predictions omit variables needed to predict future learning responses. These are different ways in which structure can arise from observational patterns. Some produce canonical quotients or completions, others identify an invariant or a family of explanations, and others select a useful representation without identifying it uniquely.

The central question is not exhausted by whether a supplied model is preserved under refinement. Nor does it require every emergence theorem to discover parts. The proposed ambition is to explain both the construction of structure and the relations between distinct constructions. A common mechanism would be valuable if it yields new consequences. A structured family may be more faithful if its divisions and directed comparisons are proved. An arbitrary common category is not enough in either case.

\paragraph{Two proposed outcomes.}
The \emph{comparative objective} asks why two procedures agree, when a meaningful translation exists only in one direction, what it forgets, and when a specified comparison cannot exist. The \emph{generative objective} asks which structural objects, relations, operations, invariants, or explanatory alternatives can be derived from specified observational material, or how learning selects among them. Comparison must do more than list separate algorithms. Generation must do more than rename supplied architecture or count fitted parameters. Both outcomes are proposed programme-level positive commitments: comparison adds an account of agreement and loss that separate constructions do not supply, while generation tests whether the theory produces or explains useful structure rather than only relating supplied models. One theorem can satisfy both if its conclusions do both jobs. Adopting this paired architecture and setting its priority remain decisions, not mathematical necessities established by this draft.

The two-sided note \cite{TwoSidedNote} motivates testing correlations, admissible realizations, specified reconstruction operations, and preservation/obstruction theorems. This is a plausible research centre, not a doctrine to which every example must conform. Its use is warranted only where input and output tests, contextual actions, and realization constraints have an independent meaning and the presentation yields a result. A one-sided, temporal, approximate, or dynamically indexed treatment remains permissible where it does better. A claim that two-sided presentation is necessary would require a separate theorem about a declared class of alternatives.

Three claims remain separate: a fundamental problem form, canonical mathematics relative to stated inputs, and a new result beyond the strongest prior work. None entails either of the others. A universal property can be genuine without yielding an executable native object. A useful learned representation can be genuine without being uniquely forced by evidence. A single impossibility result can be important without establishing that every unification must fail.

\paragraph{Independent criticism.}
The programme should be abandoned or revised if its advertised common structure contributes no theorem beyond its components, if its proposed generative novelty is already supplied in the input, or if its assumptions remove all difficult examples. Its present weakness is the absence of a demonstrated new bridge tying observationally warranted structure to the contrasting semantics of completion, prediction, native realization, and learning. The revision widens the admissible forms of a successful answer while making that substantive burden explicit.

No theorem here depends on the uncertain identification or dating of an earlier problem note. The inspected two-sided note is used for its actual mathematical content, not as evidence of an unverified authorship date or as a substitute for the current question.

\section{What the existing work already establishes}
\label{sec:ledger}

The following ledger distinguishes the wider line of work from its recent finite stochastic testbed. ``Established here'' means proved or reviewed within the stated scope of the supplied notes, not formally verified or certified novel. Recovered exploratory claims have weaker standing than the six source-reviewed reports; their prerequisites must be rechecked before reuse.

\subsection{The wider mathematical baseline}

\begin{description}[style=nextline,leftmargin=0pt]
\item[Generated observable structure]
A supplied free/forgetful doctrine and response algebra yield a free extension of preparation profiles followed by image factorization. This includes the ordinary response quotient, meet completion, linear span, and convex hull. The generated image is initial among observable realizations and terminal among generated realizations in the appropriate categories. The Boolean/FCA comparison, preservation of generator relations, test-infomorphism comparisons, and explicit observation-change obstructions are already present \cite{Core,Recovered}.

\item[Predictive reconstruction]
For a fixed initial law and specified public adaptive protocols, histories can be identified by full future experimental laws. Filtering, reachable restriction, and observable minimization are distinct operations. Filtering has the stated adjoint direction in its supplied category. A finite hidden generator can have an infinite predictive quotient; hiding observations can require filtering afresh rather than induce a map on the old predictive states. Action translation and output hiding have different comparison laws \cite{Recovered}.

\item[Linear and positive realization]
Finite-rank Hankel reconstruction and its reachable/observable properties coexist with a nontrivial choice of invariant positive cone. Generated and test-positive cones give boundaries; polyhedral and simplicial cone criteria connect to classical finite and rank-sized positive realization. Linear exactness need not preserve the generated cone. Correlated preparations can defeat separately reconstructed local tensor spaces. These are substantive results and classical antecedents, not missing generic decompositions \cite{Recovered,Positive}.

\item[Two-sided cylinder normalization]
The completed cylinder and cross-doctrine work supplies coherent normalization spans in its adequate cylinder-factorization settings. It treats actual joint preparations/tests, fixed-port wiring, signed context variance, native comparison spans, and local exactness. The family-level and substitution coherence statements are already proved under their hypotheses. Some simultaneous changes require a span between apices rather than an ordinary map of spans. Descent from a free Isbell envelope may be only oplax and need not preserve the particular mixed pullback. Wire-type changes and trace need additional structure. This proposal must dovetail with that machinery, not replace it by a new untyped square \cite{Recovered,Garner}.
\end{description}

\subsection{The recent finite native testbed}

\begin{enumerate}
\item \textbf{Common generated images and native boundaries.} The core audit isolates the finite-simplex boundary and shows why an abstract convex image does not determine its latent stochastic cost. Distinct embeddings of an abstract square can have different minimum latent cardinalities. Separation of observed rows need not separate their algebraic completion \cite{Core}.
\item \textbf{Retaining a chosen model.} The refinement note separates coarse trace preservation from a strong labelled stochastic comparison to a particular chosen realization. Bounded synthesis reduces to existential real algebra. The two-stage class reduces exactly to Blackwell garbling and constrained nonnegative rank. Its partition family gives exact feasibility and costs, including an all-budget obstruction to retaining a late-randomness model \cite{Refinement}.
\item \textbf{A common native fragment.} Normalized Boolean and nonnegative-real factorizations share a finite two-stage construction. Independent support and convex criteria decide the specified unbounded feasibility problem; explicit formulas transport witnesses and give resource inequalities. Support need not reflect quantitative feasibility or optimal size \cite{Common}.
\item \textbf{Joint witnesses and reconstruction.} Fixed local factors, occurrence labels, and statewise continuation rows give a finite tensor polytope. Feasibility constructs an actual native witness, with a dual obstruction and a sharp constrained-rank resource law. The split reconstruction is initial in a \emph{fixed tensor fibre}, not a canonical optimizer across all tensors. The note also gives direct one-step transition-witness assembly without presupposing a global readout, and a simultaneous-cover test for committed deterministic intermediate diagrams \cite{Joint}.
\item \textbf{Fixed budgets.} The gated polynomial encoding is exactly equivalent to a fixed native state budget and its displayed quartic is coercive. Global zero, positive certified global minimum, stationary point, and small numerical residual have distinct meanings. The compositional gradient theorem supplies a field-assembly identity, not a global solver \cite{Budget}.
\item \textbf{Changing interfaces and optimization.} Lawful native transport gives
\[
 \kappa(D')\leq \min_{Z\in\calP(D)}\kappa_{D'}(F_*Z)\leq\kappa(D).
\]
Both inequalities can be strict for different reasons. Inherited value functions compose by classical infimization, while independently defined native target costs need not agree with them. Canonical optimum selection, smooth slack lifting, off-zero loss preservation, and forward gradient transport can fail, including in the actual quartic rather than an unrelated surrogate \cite{Interface}.
\end{enumerate}

These six reports mostly explore one finite stochastic architecture. Their accumulation is not evidence that arbitrary multistage realization, empirical discovery, or general phenomenology is finished. Conversely, the wider cylinder, Boolean, predictive, and positive results prevent treating the line of work as merely that architecture.

\subsection{What is still not supplied}

No current result derives an unrestricted scientific vocabulary from an arbitrary experience stream. None makes finite noisy data a complete exact behavioral specification, infers hidden-state intervention access from public traces, or promises a canonical minimum native model. The joint theorem does not establish arbitrary multistage construction with stochastic intermediate links, adaptive intervention, and independently chosen assembly routes. Formal compact or Isbell completions do not by themselves supply obtainable native systems. These are boundaries to investigate or retain, not defects to conceal.

\section{The relative data of a phenomenological problem}
\label{sec:data}

The following is proposed working language, not a definition imposed on every instance in advance. Its adequacy is itself tested by the examples.

\begin{definition}[Evidence and experimental presentation]
At a stage $t$, a problem presentation distinguishes the following inputs.
\begin{enumerate}
\item A typed vocabulary $\Sigma_t$ and a \emph{currently admissible experimental repertoire} $\calE_t$. The repertoire specifies preparations, observations, sequential and adaptive protocols, parallel use, training interventions when allowed, wiring, resets, and their domains of definition. A syntactically expressible protocol need not be experimentally available.
\item \emph{Available evidence} $\mathsf E_t$: records of which protocols were performed, their outcomes, and relevant sampling, measurement, timing, and selection information. For deterministic exact measurements it may constrain specific response values; for samples it constrains laws only through a justified statistical model. The performed protocols form $A_t\subseteq\calE_t$.
\item An outcome semantics $\calB_t$ and a class $\mathsf H_t$ of candidate exact response laws on $\calE_t$. A supplied ideal law $H_t\in\mathsf H_t$, when available, describes all currently admissible experiments. It is stronger input than the available records. Alternatively, an uncertainty family $\calU(\mathsf E_t)\subseteq\mathsf H_t$ describes laws compatible with the evidence under stated deterministic or statistical assumptions.
\item A separate hypothetical repertoire $\widehat{\calE}_t$ and candidate laws $\widehat H_t$ for additional, idealized, counterfactual, or inaccessible experiments, when such an enlargement is under discussion. If $\calE_t\hookrightarrow\widehat{\calE}_t$ is given, restriction relates the two. The larger law is not inferred merely because an extension exists. Without a faithful interpretation of old experiments, even this inclusion must be replaced by a declared translation.
\item Native model classes $\calN_{\Sigma_t,d}$, their observational interpretations, and the inference or realization doctrine $d$. Algebraic closure, an architecture class, a grammar, positivity, and prior assumptions are visible inputs, not automatically discovered structure.
\item Commitments $C$ to retain, such as an endpoint realization, comparison direction, mechanism, intervention response, or chosen decomposition; and separate accounts of observation, computation, and implementation resources.
\end{enumerate}
\end{definition}

An exact law is an idealization of observations, not necessarily a law of physically impossible experiments. Conversely, a hypothetical extension of the experiment language can posit responses even when complete knowledge of all currently admissible experiments would not determine them. These are distinct uses of ``ideal''. Each theorem must say which it assumes. Statistical uncertainty about one law is also different from the availability of several native models with exactly the same law.

For a fixed stage, write $\Obs_{\calE}(X)$ for a native model's law on the admissible repertoire. The compatible realization collection is schematically
\begin{equation}
 \Real_{d,C}(\mathfrak D)
 =\{X\in\calN_{\Sigma,d}:\Obs_{\calE}(X)\in\calU(\mathsf E),\ C(X)\}.
 \label{eq:consistent}
\end{equation}
When a full exact law is supplied, use $\calU(\mathsf E)=\{H\}$ as an additional exact-input assumption, not as a statistical consequence of a finite record. Morphisms must have a concrete meaning: deterministic quotients, stochastic intertwiners, algebra homomorphisms, or specified comparison spans. The collection itself is bookkeeping, not a research result.

\paragraph{Three changes with different effects.}
More evidence from the same repertoire can shrink an exact consistency set; a succession of ordinary confidence sets need not be nested, so nesting requires a construction and a coverage guarantee. Enlarging the admissible repertoire can refine observational equivalence even before new outcomes have been recorded. Enlarging the hypothetical law domain can change a theoretical target without changing available evidence or experimental access at all. Changing the native doctrine can alter realizability without changing any response law. A proposed functor or information order must specify which of these changes it represents.

For later resource formulas, $E$ denotes the explicitly chosen response interface and $\calU_E$ its compatible response set. It is not a synonym for the raw record $\mathsf E$, nor necessarily the entire currently admissible repertoire. A revised vocabulary must interpret the old protocols, and a new primitive needs an executable procedure or an explicitly conjectural status.

\subsection{Forced structure and chosen structure}

A structural predicate is \emph{forced relative to $\mathfrak D$} if it holds in every admissible consistent realization, interpreted up to the declared comparison or equivalence. It is refuted if none satisfy it and undetermined if both possibilities occur. This trichotomy presupposes a nonempty consistent class; an empty class is model failure, not vacuous evidence for every claim. A basis, positive cone, optimizer, coupling, decomposition, regularizer, or prior can be a useful choice without being forced.

A positive identification claim must name the invariant being identified. Isomorphic observable quotients, equivalent experimental predictions, equal state minima, the same retained transition mechanism, and isomorphic native models are different answers. A same-source choice of one realization does not prove that another route produces exactly the same output, or even an object in the same retained-witness fibre.

\subsection{Four quantifiers that cannot be interchanged}

For an observation change, one can ask for:
\begin{enumerate}
\item an abstract reconstruction and its comparison;
\item existence of \emph{some} native realization of the revised behavior;
\item a certificate relating \emph{specified} endpoint realizations;
\item a lift for \emph{every} admissible target realization, possibly within every budget.
\end{enumerate}
The fourth is much stronger than the second or third and is not the default demand. Useful forgetting need not lift every coarse choice. A proposed sufficiency theorem must state which quantifier it proves, rather than rejecting useful abstractions because unrestricted fixed-endpoint lifting fails.

\subsection{Resource semantics}

Use a declared ordered resource set, for example
\[
 \Cost(X)=(\text{latent states},\text{description size},\text{execution cost}),
\]
and keep experimental cost separately as a budget on protocols and samples. Addition or composition of costs requires a model; it is not inferred from category notation. State count is not automatically bits, table size, energy, time, or sampling complexity. External clocks, shared randomness, hidden correlations, preprocessing, and oracle calls must be charged or explicitly provided.

For a scalar native budget $b$ and tolerance $\varepsilon$, a basic profile is
\begin{equation}
 \mathscr B_{d,C}(\mathfrak D)
 =\{(b,\varepsilon):\exists X\in\calN_{\Sigma,d},\ C(X),\
 \Cost(X)\leq b,\ \exists H\in\calU_E:\ d_E(\Obs_E X,H)\leq\varepsilon\}.
 \label{eq:budget-profile}
\end{equation}
Here the metric or discrepancy and its relation to confidence must be specified. The displayed existential quantifier expresses \emph{some} compatible response; robust conclusions valid for \emph{every} member require a different quantifier. Exact feasibility, approximate feasibility, and robust feasibility must not share one symbol without this distinction. Retaining a resource spectrum or a set of witnesses is often appropriate when choosing an optimum coherently is impossible.

\section{Where structure can actually emerge}
\label{sec:emergence}

A useful theorem must identify structural output that was not installed in the input. Start, in a simple two-sided case, with a response pairing
\begin{equation}
 H:P\times T\longrightarrow\Omega,
 \label{eq:pairing}
\end{equation}
where $P$ and $T$ name available presentations or preparations and observations or tests. They need not describe an interactive system: a static incidence relation or a family of measured features is already an instance. Further composition or context data are included only when the observations support them or the experimental setup supplies them.

\subsection{Distinctions and states}

Equality of response profiles gives the observational relation
\[
 p\sim_Tp'\quad\Longleftrightarrow\quad H(p,t)=H(p',t)
 \text{ for every }t\in T.
\]
Its quotient contains distinctions made by these tests. In a temporal setting, take preparations to be accessible histories and tests to be continuations; equality of complete future laws generates predictive states. This exact construction assumes full behavioral semantics, an adequate exact oracle, or a justified ideal limit; a finite observed record alone does not supply those laws. The source of the state space is then a regularity of observed behavior rather than a chosen hidden carrier.

This is an identification relative to a language of tests. It does \emph{not} force every compatible native model to have literally identical underlying states for equivalent histories. Adding tests can split observational classes, while adding evidence within a fixed language can shrink a set of explanatory models. Those are different operations with different variance. A theorem must not conflate them.

\subsection{Relations and completed objects}

Incidence patterns can generate closed concepts or implications; linear dependence can generate a dimension and equations; admissible mixing can generate convex structure. A completion can introduce objects that were not individually observed. Its warrant may be mathematical closure under a supplied principle, rather than evidence that each formal object is an executable preparation.

The existing generated-image theorem already explains several such constructions. The new task is not to repeat it with more examples. It is to explain its relation to structures generated by continuation, latent factorization, and possible decomposition, including cases where their universal properties or native meanings are different.

\subsection{Operations induced by observed composition}

Suppose experimental contexts act on preparations and tests with an observed compatibility law such as
\begin{equation}
 H(c_*p,t)=H(p,c^*t).
 \label{eq:context-duality}
\end{equation}
One can ask whether $c_*$ descends to the emerged distinctions, whether its composites induce relations among the new operations, and whether those relations survive a chosen completion. The experimental context is supplied; its operation on a newly constructed state space and the relations it satisfies can be derived. Conversely, if no context action or compositional evidence is available, an invented multiplication on the state set must be labelled as additional structure.

A more ambitious question is whether patterns of contextual transformations determine a smaller operation language or reveal reusable action families. This is a discovery problem relative to the available contexts and admissible hypothesis language. It is not permission to claim all scientific operations follow from an unstructured sequence of outcomes.

\subsection{Latent carriers and decompositions}

A factorization pattern may introduce intermediate objects through which observed relations or laws pass. These objects can explain response dependencies or offer a compact implementation. Unlike a quotient or a closure, a smallest factorization need not be unique or characterized by the same universal property. Different factorization witnesses can carry different conditional mechanisms while agreeing on all present observations.

A decomposition is stronger than merely factoring a matrix through an arbitrary carrier. It posits reusable parts, interfaces, and an assembly operation. The structural question is whether patterns in the observations and available contexts warrant such parts and their reuse. A central flagship test should construct candidate interfaces or factorizations without taking the desired component split as input. Its success is not a necessary condition for every other form of emergence. A declared class of admissible decompositions is allowed; supplying the actual factors, interface, and global witness and then proving they compose is a different, downstream question.

\subsection{A concrete supplied and derived ledger}

For every example and theorem, list:
\begin{enumerate}
\item the raw observational data and experimentally available manipulations;
\item the comparison of outcomes and inference principles supplied;
\item the structural objects, operations, relations, and interfaces constructed;
\item the invariant or universal property that warrants each construction;
\item the unresolved choices and the additional observations that could distinguish them;
\item the native objects or operations that the construction does and does not realize.
\end{enumerate}
This ledger should show genuine emergence even in exact, non-algorithmic mathematics. It also prevents a supplied latent category from quietly carrying all the structure the theorem claims to discover.

\section{The comparative and generative targets}
\label{sec:central}

\begin{target}[A substantive comparative programme]
\label{target:central}
For a stated substantial class of observational problems, compare incidence/closure, temporal/predictive, and native-factorization constructions at their correct types. Establish input-level criteria explaining agreement, directed comparison, information or resource loss, and failure of a specified comparison. Connect the criteria to what each construction actually produces and supports. Prove at least one new positive bridge across unlike mechanisms beyond the existing generated-image and cylinder-normalization theorems, and a sharp boundary showing where the bridge ceases to apply.
\end{target}

This objective follows Stacey's method at the level of research design: expose the choices, derive comparisons from their change, determine their exact properties, then test stronger claims independently. A common two-sided presentation may support this work. It does not itself answer why a comparison preserves meaning, why it loses information, or whether a stronger comparison exists.

\subsection{The precise comparative precedent}
\label{sec:stacey}

Stacey's PDF Definition 3.1 uses a faithful $j:\mathcal T\to\mathcal U$, an underlying $X$, input/output hom-subfunctors $I,O$, and the requirement that every input--output composite underlie a test arrow. Definitions 3.4, 3.6 and 3.8 add trials, forcing functors, and forced objects \cite{Stacey}.

Its comparison hypotheses matter:
\begin{itemize}
\item Stronger forcing gives an inclusion (Proposition 4.3). The operators of Definition 4.5 need not satisfy both adjunction inequalities; Proposition 4.7 supplies conditional results, and Section 4.2 supplies a counterexample.
\item For a full test inclusion, restricted forcing and $\mathcal U$-adequacy (Definition 4.10) yield $\mathrm{Ext}^{-}\dashv\mathrm{Res}\dashv\mathrm{Ext}^{+}$ (Corollary 4.16). Adequacy detects each inadmissible underlying arrow by sandwiching it with smaller-test arrows.
\item Base change must respect tests (Proposition 4.18); adjunction transfer and restriction to forced objects have further hypotheses (Corollary 4.19; Lemmas 4.20--4.22; Propositions 4.25--4.26).
\end{itemize}
Sections 6--7 separately exclude further adjoints; Section 8 uses invariants for non-equivalence. Thus unavailable proofs of stronger properties are not proofs of their failure. Stacey constructs structure as well as comparisons; it is no mere organizational metaphor. Its smooth-space hypotheses do not automatically supply empirical identification or native realization. Numbering refers to the inspected PDF.

\subsection{The semantics of a comparison}

For each pair of methods, specify their input regimes, output categories, intended observational or structural semantics, and an independently understood common fragment. A proposed translation must state what it preserves on that fragment. The possible findings include:
\begin{enumerate}
\item an equivalence or an isomorphism of constructed outputs, with the required preservation hypotheses;
\item a directed functor, reduction, natural comparison, or typed span, without an unjustified inverse;
\item a quantified loss, such as distinctions forgotten, a completion no longer preserved, higher native resource cost, or controlled approximation error;
\item an obstruction to the specified comparison class, possibly despite equal present predictions;
\item absence of an adequate shared interpretation, in which case a comparison is not yet well-posed.
\end{enumerate}
A failed candidate functor does not establish that no useful comparison exists. A failure of equivalence does not imply absence of directed functors. Two incomparable elements of a chosen information preorder need not define incomparable categories. An inability theorem must quantify over a specified family of constructions or translations, with its allowed semantics, common-fragment constraints, resources, and desired preservation properties. ``No categorical unification'' is not a theorem statement.

The original note's residual-language inclusion into intersection closure and its rank-three versus nonnegative-rank-four example are controls \cite{TwoSidedNote}. They already demonstrate a positive bridge and a native obstruction. Reproducing them clarifies the problem but does not meet the new-result obligation. The same applies to a preservation theorem obtained immediately by essential uniqueness of a designated factorization.

\subsection{The distinct generative objective}

\begin{target}[Generation with a structural consequence]
\label{target:generative}
From an explicitly specified evidence or ideal-law presentation, with a declared repertoire and inference doctrine, construct a response-relevant structural output not supplied as that output in advance. Prove a positive theorem that explains its generation or learning-driven selection, with an independent criterion for its warrant and a sharp ambiguity or failure boundary. Relate the generated structure to at least one unlike reconstruction mechanism through the comparative programme, without requiring the two carriers to coincide.
\end{target}

The structural output can be objects or states, operations or relations, a geometric or algebraic invariant, a sufficient dynamic state, a hypothesis language, a family of latent explanations, or a candidate interface or decomposition. An informative structure already present in a broad search space can be selected by learning; the theorem must explain the selection and its observational consequence rather than claim that the whole search space was discovered. No one output form is mandatory. A useful exact construction remains emergence even when it is neither algorithmic discovery nor a decomposition.

\paragraph{Nonvacuity requirements.}
A claimed generative result must meet all of the following.
\begin{enumerate}
\item \textbf{Specified output and use.} Name the invariant, structure, or equivalence class and the predictions, compositions, interventions, or learning responses it makes possible. Define the equivalence under which it is identified.
\item \textbf{Visible supplied structure.} Inventory the experimental vocabulary, inference principles, architecture, dynamics, initialization, and optimization choices. Separate generation from the mere exposure, fitting, or relabelling of a supplied structure.
\item \textbf{Independent input criterion.} State assumptions as checks, promises, invariants, or structural hypotheses on the input. Do not restate ``the desired realization exists'' or build its unknown solution into a certificate. A checkable witness is useful only with a theorem connecting independently obtained witness data to the result; general search over every possible solution is not by itself the new criterion.
\item \textbf{Positive content.} Give a construction and proof on a nontrivial family, not only an absence theorem or an easy example. The family must retain a declared source of genuine inference, selection, or synthesis difficulty.
\item \textbf{Counterfactual dependence and robustness.} Show which changes in observations or learning conditions change the structural result, and which admissible relabellings or reparameterizations do not. A structure returned unchanged for every admissible input needs an independent explanation of its evidential role.
\item \textbf{Warrant and limits.} Distinguish what is identified, doctrine-relative, statistically supported, or selected by an objective; exhibit observationally compatible alternatives where they remain. State why the new theorem is not already a cited predecessor in changed notation.
\end{enumerate}

A structure can be new to an instance while its construction theorem is classical. Conversely, a new theorem can explain a familiar kind of structure. These two novelty claims must be audited separately.

These criteria do not insist on discovering a vocabulary from an arbitrary unstructured stream. They do insist that the theorem explains a substantive dependence between the permitted input and the output. A universal free construction may be mathematically meaningful; to count as the new generative result it must have an additional structural consequence beyond the existing generated-image baseline.

\subsection{A demanding flagship rather than the sole gate}
\label{sec:flagship}

Retain the earlier challenge: start with response patterns and a finite presentation of available contextual transformations, without naming the desired components. Construct or obstruct a reusable factor presentation. In a deterministic case this may involve congruences or a transformation algebra. A stochastic version must retain positivity, correlated observations, nonunique factors, and the distinction between predictive and hidden state. Its output can be a family of possible decompositions with a stated equivalence and discriminating probes.

Rule out identity factorizations and one-part decompositions in the first finite test. Require at least two proper observable parts, each carrying distinctions used by actual tests or contexts, and an assembly reconstructing the claimed joint responses and operation action. Product-style examples can use proper projections, joint separation, and verified action/readout equations. Overlapping interfaces require their own nontriviality criterion. Projections, context assignments, and assembly witnesses are outputs, not supplied answers. Completeness assumptions are necessary when extending a finite verification to untested protocols.

This flagship remains open and demanding. Meeting another genuine generative target does not mark it completed. Failure of this particular form does not invalidate a theorem generating other structure. Likewise, an impossibility theorem for a proposed common mechanism is valuable but does not automatically discharge either positive target. Replacing a positive target by a classification-only programme requires an explicit change of the provisional contract.

\subsection{Strong emergence precedents and their limits}

Computational mechanics constructs predictive states and transitions from process laws relative to a supplied temporal and predictive framework \cite{ComputationalMechanics}. FCA derives concepts and lattice operations from incidence. CALF constructs image hypotheses and compatible algebra structure under its stated categorical and learning assumptions \cite{FCA,CALF}. These are substantive precedents, not counterexamples to emergence because no interface is discovered.

Contextual grammar learning, structural-form learning, and program-library learning move other parts of the supplied/derived boundary \cite{ContextualGrammar,StructuralForm,DreamCoder}. Their input languages, hypothesis spaces, queries, and inductive biases must remain visible. Causal macro-variable construction requires its stated causal and intervention assumptions \cite{MacroVariables}; nonidentifiability results are equally model-relative \cite{Disentanglement}. The research obligation is to explain relations and obtain an additional result, not to relabel all these works as instances and claim their achievements again.

\section{Learning dynamics as a candidate generative test}
\label{sec:learning}

Learning dynamics deserves a prominent place among the candidate tests, alongside static incidence and exact reconstruction. It exposes distinctions invisible in current predictions and permits controlled changes of data, update rules and representations. These features make it a plausible test of how structure is generated or selected, rather than just a solver implementation. The counterargument is substantial: a supplied architecture and optimizer can carry the apparent structure in advance, a sufficient learning-state reconstruction can reduce to established realization theory, and results on a tractable training task may not transfer. Its adoption as the first generative test therefore remains open. A training run by itself proves neither emergence, identification, nor novelty.

Let $\Theta$ be a supplied parameter or implementation space with a response map $h$ and a family of declared learning updates $U_a$ or flows, indexed by allowed training inputs or interventions $a$. Present predictions identify parameters with the same $h(\theta)$. That quotient need not support learning: two currently indistinguishable parameters may respond differently to the same future training protocol. For a declared probe family $\Pi$, define the response equivalence
\[
 \theta\sim_\Pi\theta'
 \quad\Longleftrightarrow\quad
 \text{all permitted learning-and-observation protocols in }\Pi
 \text{ have the same response law from }\theta,\theta'.
\]
Reconstructing a sufficient state for a supplied learning system is different from explaining how a new representation forms or is selected during training. The first is an observational reconstruction problem; the second needs a theorem or experiment about the formation/selection process itself. The exact target below primarily addresses the first; the optional empirical branch can test the second. Neither achievement should be reported as the other.

This is a target semantics. Merely forming its quotient or retaining all of $\Theta$ is not the new theorem. The hard question is whether a finite or otherwise tractable response-defined structural presentation closes under the learning dynamics, which information it necessarily retains, and how it changes with evidence, access, or learning rule.

\subsection{An elementary falsifying control}

For the scalar predictor $w=uv$ and Euclidean gradient flow of a $C^2$ loss $\ell(w)$, on a common interval of existence,
\[
 \dot u=-v\ell'(w),\qquad \dot v=-u\ell'(w),\qquad
 \dot w=-(u^2+v^2)\ell'(w).
\]
Thus $(u,v)=(1,1)$ and $(2,1/2)$ have the same current predictor $w=1$, but different initial predictor velocities whenever $\ell'(1)\ne0$. With $s=u^2+v^2$,
\[
 \dot w=-s\ell'(w),\qquad \dot s=-4w\ell'(w),\qquad s\ge 2|w|.
\]
The pair $(w,s)$ closes this specified flow. This exact calculation is an elementary control, not a new theorem, a claim of global minimality for every probe family, or evidence from a training experiment. If no training intervention or future training-response test is admissible, the larger observational distinction is not thereby accessible. If the chosen loss has $\ell'(1)=0$, the displayed first-order discriminator vanishes. For a sufficiently rich admissible family of losses and derivative or short-time observations, the additional distinction can be operationally tested.

The example shows why an interface or parts cannot be the sole generative output: an observationally relevant learning-state variable can matter without decomposing the system. It also shows the danger of confusing the function represented with the representation's induced learning dynamics.

\begin{target}[Response-defined structure under learning]
\label{target:learning}
For a declared, nontrivial family of learning systems and experiments, derive a tractable presentation of the structural variables or relations needed for future learning responses. Prove closure or a controlled approximate closure, together with a stated sufficiency, minimality, or unavoidable-ambiguity result relative to the probe family. Prove a positive comparison with an unlike reconstruction mechanism or between two materially different learning regimes, and identify a sharp obstruction to stronger invariance or compression. The construction must do more than retain all parameters or fit a prespecified latent architecture.
\end{target}

A candidate route is a response-generated algebra of observables closed under declared update pullbacks or flow derivatives, followed by a finite-generation or finite-dimensionality criterion. Its closure may be infinite, its invariant may be nonidentifiable from samples, or its construction may simply reduce to known realization theory. Nonlinear realization, observable-algebra closure, and Koopman-style approaches are explicit reduction targets before any novelty claim. All three possibilities must be tested before claiming novelty. A theorem about a useful finitely generated class and a boundary to closure would be substantive if its extra content survives those reductions.

\subsection{Controls for a learning claim}

A learning result must name the object that emerges, the phase or limit at which it is claimed, and its operational use. Track the induced behavior and an invariant structural measure alongside the training loss. Compare with initialization, random features or a fixed-representation baseline where relevant, and alternative admissible dynamics. Control the architecture, optimizer, regularizer, data ordering, and access assumptions that could already install or select the answer. Selection by these choices can be a legitimate result, but not evidence that observations uniquely force that structure.

A pure coordinate change with the learning vector field and all observations transported appropriately must preserve the physical prediction; ordinary Euclidean gradient descent redefined after a nonlinear coordinate change is generally a different learning system. The theorem must separate these cases. Equivalent predictors can induce inequivalent training responses, and equal training loss can conceal different generalization or intervention behavior. Apparent emergence that disappears under an admissible gauge transformation, a held-out protocol, or an architecture control fails the stated invariant claim.

When finite records are used, specify which learning responses were observed, the sampling assumptions, confidence or convergence interpretation, and the additional experiments needed to distinguish compatible alternatives. A fitted low-dimensional description is evidence for a candidate, not a proof of exact closure or globally minimal structure. Any causal claim about training must name the actual interventions, the variables manipulated, and the assumptions connecting them to the causal question.

\subsection{An empirical learning test if that branch is selected}

A controlled neural-learning example can test discovery rather than only exact reconstruction. Nanda and collaborators already infer a Fourier-based mechanism for modular addition and use intervention-based progress measures; their discussion identifies manual, task-specific analysis and unresolved ex ante transition prediction as limits \cite{Grokking}. Reproducing that mechanism would be a control, not new discovery.

A demanding proposed test would infer explanatory coordinates using only information available before a declared training cutoff, hiding the final checkpoint and evaluation interventions during discovery. Assess predictions on unseen runs and interventions against ordinary loss curves, fixed representations, and relevant mechanistic baselines. Predeclare the prediction, intervention, resource, or sample-use advantage being tested. Include behaviorally similar systems with different internal mechanisms and controls for architecture and optimizer choice. A post hoc metric fitted after seeing the transition fails the claimed prospective test.

This is one possible empirical instantiation, not an added commitment to run neural experiments or a prerequisite for every pure theorem. Its discovery procedure and mathematical warrant would need their own strongest-source audit. A successful benchmark alone would not establish the comparative programme or a general causal discovery theorem.

\section{Refinement laws as a component of the programme}
\label{sec:refinement-component}

After a structure has emerged, one can ask what survives a richer test language, a changed doctrine, a new connection, or a replacement mechanism. This is refinement-law territory, with substantial prior art. The existing programme already contains strong comparison results, and abstract interpretation supplies soundness, completeness, and constructive domain repair; refinement mappings can require history and prophecy variables; interface theories give compositional substitutability under explicit compatibility conditions \cite{AbstractCompletion,RefinementMappings,InterfaceAutomata}.

Generalized strong preservation already connects observation languages, abstract domains, behavioral equivalences, and complete-domain construction; it is an especially strong common-theory control under its stated hypotheses \cite{StrongPreservation}. Predictive sufficiency can itself be a constructive emergence principle. The issue here is the narrower preservation question that starts from an already selected representation, not the word sufficiency.

The proposed refinement component should therefore answer a specific residual: how do the newly generated objects, operations, and candidate interfaces compare, and when can those comparisons be realized natively with the stated resources? The finite sufficiency test below is a demanding positive special case. Noninvertible comparisons, controlled forgetting, and set-valued explanations remain first-class outcomes. Universal same-budget lifting is not the definition of useful phenomenology.

\section{A multistage test for emerged interfaces}
\label{sec:multistage}

The recent native work gives a stringent downstream test of an emerged interface. This target remains a substantive positive obligation of that flagship, not a mandatory definition of all emergence. It does not replace Target~\ref{target:generative}; it asks whether such an interface can actually be used and revised.

\subsection{Input and requested output}

Take finite controlled systems with joint output and next-state kernels
\[
 M^i_{a,y}:S_i\longrightarrow\RR_{\geq0}^{S_i},
 \qquad \sum_y M^i_{a,y}\one=\one,
\]
finite supplied preparations, an explicitly declared finite operation/intervention repertoire, and a diagram of partial local observations and retained models. Include the prescribed deterministic or stochastic comparison maps with their directions. Require sequential state reuse through at least two successive nonterminal transitions: a state reached by the first operation can be acted on and observed by the next. When feedback is claimed, additionally allow a later action to depend on an earlier observed outcome. At least one intervention must replace or control a specified mechanism, rather than merely relabel an output. The permissible intervention semantics and access to it must be supplied; passive observations do not create that access.

A candidate joint implementation has its own finite carrier $V$, kernels $X_{a,z}$, preparations, and retained links $L_i:V\to\Delta S_i$ where stochastic links are intended. For a declared observation map $q_i$, a required retention equation has the concrete form
\begin{equation}
 \left(\sum_{z:q_i(z)=y}X_{a,z}\right)L_i
 =L_iM^i_{a,y},\qquad \pi_XL_i=\pi_i.
 \label{eq:multistage-retention}
\end{equation}
A stochastic output channel uses the corresponding weighted sum. Analogous equations are required for each declared intervention, with the exact operators it changes. These are statewise conditions, stronger than equality of the one initial trace law. More general native classes may require a different typed formulation; Equation~\eqref{eq:multistage-retention} is the first concrete target, not a definition of all phenomenology.

The input must not already include a full joint witness. It may include incompatible or only partially determined local pieces, and previously chosen intermediates can matter. The requested output is a construction or obstruction for a joint native realization and for at least two specified assembly routes which begin from those pieces. Rebracketing a single supplied full tensor is a regression case, not the new theorem.

\begin{target}[Native use of an emerged interface]
\label{target:multistage}
Find an explicitly described class of the preceding inputs, closed under a useful collection of sequential operations and mechanism interventions, with an independent finite sufficiency/compatibility criterion which:
\begin{enumerate}
\item constructs joint native witnesses from the local presentation, with no assumed global coupling or unknown global readout supplied as a premise;
\item retains the specified stochastic links and operational comparisons, or gives an exact criterion distinguishing when only weaker public behavior can be retained;
\item compares two independently built legal assembly routes by actual native morphisms or a justified typed replacement, including any genuinely additional compatibility condition;
\item proves a nontrivial native resource bound and a matching lower bound on a parametrized family whose difficulty is not eliminated by a free clock, a full-history state, or a supplied witness;
\item identifies how a finite observational sufficiency certificate is used in the construction and how it changes when a test, operation, or intervention is added or hidden.
\end{enumerate}
\end{target}

A theorem for all finite systems is not presumed. Candidate restrictions include a finite continuation basis, a controlled overlap pattern, a specified cone class, bounded interface width, or a verifiable stability/separation condition. Their meaning and testing cost must be stated. A chosen class must still contain genuine stochastic intermediate links, unknown joint choices, sequential reuse, and a nontrivial intervention or its precisely diagnosed failure. A restriction that reduces every case to one already solved row polytope fails the milestone.

\subsection{What counts as an independent criterion}

Acceptable possibilities include a finite relation/closure calculation that yields the necessary native witness, a dual certificate with a construction theorem, or a finite positivity and invariant-subspace condition whose sufficiency is proved for the declared class. A finite cone or rank certificate can be useful when the generators and bounds are supplied; existentially asking for an arbitrary invariant cone is not an effective unbounded criterion by itself. A criterion stated as ``there is a compatible joint realization'' or ``the operation preserves the relevant class'' must be replaced by an input-level test.

The proof must reduce standard cases to their established explanations: ordinary marginal extension, conditional products, supported coupling, Markov semipullbacks, multicausal gluing, and bounded real-algebraic synthesis. The additional theorem must identify exactly what survives after prescribed stochastic anchors, finite memory, or alternate assembly routes are imposed. Merely adding the word ``mechanism'' to a fixed conditional-marginal constraint is not a new result \cite{Joint,Refinement,Multicausal}.

\subsection{Negative outcome and repair}

If a proposed condition is insufficient, construct a minimal or structurally transparent counterexample with exact verification. Then state the failed conclusion and test a repair: enlarging the retained interface, retaining a set of joint witnesses rather than one optimum, allowing posterior states, adding a necessary probe, changing the comparison class, or explicitly relaxing a retention requirement. The repair must preserve a useful operational claim and have a proved positive theorem with resource accounting.

A counterexample can refute the proposed unrestricted statement. It cannot satisfy all the positive clauses of Target~\ref{target:multistage}. If no nontrivial repaired class is found, this milestone remains incomplete and the broad version of the proposal must be revised. A tiny example that isolates an obstruction is valuable; it is not a substitute for a family retaining substantial inference or synthesis difficulty.

\section{A positive sufficiency test within the refinement component}
\label{sec:finite-sufficiency}

A demanding positive instance makes the refinement component more falsifiable. It is deliberately stronger than ordinary forward forgetting, which remains useful when it fails.

Let $E\subseteq E^+$ be finite observation interfaces in a declared finite-dimensional or finite-presentation regime, with restriction $q$. Let $\mathcal G$ be specified generating operations and exposed interventions, with their allowed wiring rules. Fix a native budget cap $B$. For an exact full response specification $H$, write $\Real^+_b(H,C)$ for its native implementations with cost at most $b$, and $\Real_b(qH,qC)$ for the restricted problem, using the correctly translated commitments.

\begin{target}[Finite operational sufficiency with native lift]
\label{target:finite}
Derive finite input certificates, under stated structural assumptions, which imply all of the following for every $b\leq B$:
\begin{enumerate}
\item observations on $E$ determine the required $E^+$ predictions throughout the protocols generated by $\mathcal G$, not just on the originally measured preparations;
\item the budget feasibility statements agree,
\begin{equation}
 \Real^+_b(H,C)\ne\varnothing
 \quad\Longleftrightarrow\quad
 \Real_b(qH,qC)\ne\varnothing;
 \label{eq:budget-exactness}
\end{equation}
\item for the stated class of coarse implementations, the certificates construct same-budget native lifts and comparisons, with a composition law under the generated wiring and intervention operations;
\item a corresponding weaker certificate gives a controlled one-sided comparison, resource bound, or explicit missing-interface obligation when exactness fails.
\end{enumerate}
The hypotheses must be verifiable in the presentation, not a prior assumption of a reverse native functor, budget-preserving lifting, or coherence of the desired outputs. The theorem must say whether its native conclusion covers every coarse implementation, only specified endpoints, or only existence.
\end{target}

Equation~\eqref{eq:budget-exactness} alone is weaker than the lifting clause. For attained scalar integer costs, equality of minima is equivalent to agreement of feasibility at every budget; agreement only for $b\leq B$ can hold while both minima exceed $B$ and differ. A multidimensional resource frontier can carry more information than a scalar optimum. The purpose is to derive an operationally meaningful native lift from a finite sufficiency certificate, rather than prove equality by assuming lawful transports in both directions, which the existing interface note already does.

\subsection{A clean control lemma}

There is a simple finite prototype that should be proved and used as a control, not counted as completion. Suppose the \emph{full admissible latent-profile polytopes} $K,K'$ are supplied, and a linear observation map $q$ satisfies
\begin{equation}
 q(K)=K',\qquad
 \ker(q)\cap\spanlin(K-K)=\{0\}.
 \label{eq:polytope-iso}
\end{equation}
Finite rational vertex/facet and rank checks can certify these conditions for appropriate explicit presentations. Require the actual generating operations to satisfy their well-typed intertwining equations on sufficient generators, and require explicitly specified joint correlation carriers as well as the local polytopes. Then an affine inverse on the admissible profiles is available without requiring a stochastic inverse on the entire ambient simplex. One can ask how that inverse lifts finite latent factorizations and the declared operations while keeping the same hidden carrier.

This is stronger than injectivity on observed rows and potentially weaker than a globally invertible observation channel. It still supplies much of the needed structure in advance. The resulting conjugacy argument may be elementary. It must be compared with existing tomography, statistical sufficiency, automata minimization, and causal abstraction before any novelty claim.

\subsection{The hard bridge}

The substantive refinement extension is to derive usable certificates of the relevant full admissible structure from \emph{available observational and operational data}, instead of supplying it as $K$. In positive linear settings, the existing generated and test-positive cones suggest a candidate interval
\[
 C_-(E)\subseteq C\subseteq C_+(E).
\]
One must track the observation embedding, normalization, all continuation/intervention operators, and retained native links. Equality of suitable bounds can make a cone forced; a gap records genuine freedom. Neither local cone equality nor a basis of marginal tests determines arbitrary joint correlations. Nor does a polyhedral generated cone automatically describe every latent profile admitted by a native implementation.

A useful theorem would identify finite conditions, such as separation, positivity reflection on the relevant span, operational closure, and controlled correlation structure, which imply the native conclusion in Target~\ref{target:finite}. It must derive the native lift rather than conceal it in the conditions. Conversely, an obstruction should distinguish insufficient observation from a genuinely incompatible retained model or resource budget. At least one criterion should remain useful when the canonical predictive object is larger than the hidden implementation, and at least one should expose the failure of purely local observational adequacy after composition.

The same pattern should have an algebraic instance: preservation of observed generators is insufficient unless the relevant generated relations and operation uses are preserved. The common theorem should explain why relation preservation, predictive continuation closure, and positivity/realization conditions occupy comparable roles, without falsely identifying their native categories or adjunction directions.

\subsection{What would refute this proposed direction}

If every plausible finite certificate either encodes the entire sought native solution, reduces to an existing theorem without added content, or excludes the mandatory hard cases, then the proposed sufficiency theorem has not been obtained. A rigorous classification showing that the mechanisms cannot admit the proposed common form would be a serious theoretical outcome. It would support a structured family of theories with sharp bridges, rather than a universal reconstruction. Such a result must establish the obstruction at the stated level of generality; one easy counterexample does not settle the fundamental-status question.

\section{Comparison with the cylinder and context machinery}
\label{sec:cylinders}

The existing two-sided machinery is a mathematical dependency and a test of the proposal. Garner's Definition 2.2 requires designated orthogonal cocones and cones factoring every small cylinder; Theorem 3.3 gives the free cylinder-factorization completion and Theorem 4.1 its pseudoalgebra characterization. Theorem 5.1 distinguishes lax/colax morphisms preserving the respective classes \cite{Garner}. These universal properties and their size assumptions must not be silently identified with Stacey's forced structures or with empirical emergence. In each supported instance, supply an interpretation
\[
 J_{\mathfrak D}:\Real(\mathfrak D)\longrightarrow\operatorname{CylReal}(\mathfrak D)
\]
into the appropriate cylinder-realization setting. This map must send the \emph{actual} preparations, tests, native comparisons, and declared operations to their two-sided data. It must record which native information is forgotten. A freely completed formal object is not silently substituted for its native preimage.

For an honest cylinder realization $X$, denote the established normalization span schematically by
\begin{equation}
 \nu_{\mathfrak D}(X)=
 \left[X\xleftarrow{\mathcal M}G_X
             \xrightarrow{\mathcal E}R_{\mathfrak D}\right].
 \label{eq:normalization-span}
\end{equation}
The actual classes and variance are those of the existing theorem. A supported change $F$ has its established cross-doctrine comparison
\begin{equation}
 \beta_F:
 F(R_{\mathfrak D})\xleftarrow{\mathcal M}G_F
                \xrightarrow{\mathcal E}R_{\mathfrak D'}.
 \label{eq:cylinder-beta}
\end{equation}
These are baselines, not conjectures of this proposal.

The new theory must give a specified comparison from its interpreted native construction to these spans. Where native forward transport $u_*$ exists, the first required square is
\[
\begin{tikzcd}[column sep=large,row sep=large]
\Real(\mathfrak D) \arrow[r,"u_*"] \arrow[d,"J_{\mathfrak D}"'] &
\Real(\mathfrak D') \arrow[d,"J_{\mathfrak D'}"] \\
\operatorname{CylReal}(\mathfrak D) \arrow[r,"F"] &
\operatorname{CylReal}(\mathfrak D').
\end{tikzcd}
\]
The claim may be strict commutation, a specified natural comparison, or an appropriate span-valued version; the theorem must choose and type it. It must then prove that normalizing this square produces the existing comparison~\eqref{eq:cylinder-beta}, up to a specified isomorphism or noninvertible comparison of spans. Equivalently, the native route and the established two-sided route must be related by an actual morphism or two-cell in the supported span calculus, not by the phrase ``compatible with cylinders.''

For consecutive changes, this comparison must respect the established compositor and its unit and associativity laws. For simultaneous signed changes in preparations and tests, it must use the surviving span-between-apices construction where an ordinary interchange map is impossible. For fixed-port wiring it must compare the actual joint cylinder; for wire-type change it must supply the extra transport and account for reversed dual variance. Feedback or trace requires the original definedness and preservation conditions. Monoidality alone does not prove trace preservation.

The decisive extra content is reflection toward native semantics. The two-sided image may exist although no admissible finite native realization does. A faithful statement must expose the resulting lifting obstruction or resource cost. Likewise, a functor preserving factorization classes can transport a comparison without preserving the relevant mixed pullback exactly. The proposed sufficiency certificate should tell when native conclusions can be recovered, rather than simply reasserting the free comparison theorem.

Contextads are a possible language for the dependence of operations on context, revision, or partial undo. They become appropriate only if an explicit construction in this problem satisfies their required structure and yields a result not already given by the cylinder machinery. An institution, double category, fibration, profunctor, or operadic doctrine is similarly a candidate tool. No one of these is selected as the answer before the native criteria and comparisons are understood.

\subsection{Causal interpretation is an additional claim}

A native state, factorization, or learned variable need not represent a physically intervenable mechanism. Equality of passive or public trace laws does not identify hidden-state interventions. An intervention claim must specify its admissible manipulation, its target and scope, how it acts on the law or native model, and the assumptions supporting transport to the proposed causal interpretation. Conditioning, output hiding, and replacing a mechanism are different operations. Hypothetical intervention laws belong in the enlarged-law input unless their experimental accessibility is established.

A causal acceptance test compares the stated intervention responses, not only observational fit. A mandatory falsifying control is a pair of admissible models with the same available observational evidence but different permitted or hypothetical intervention laws. If the interventions are inaccessible, the appropriate conclusion is an identified set or conditional statement, not a causal identification theorem. If they are accessible, the acquisition procedure must say which experiments distinguish the models and at what cost.

\section{Acquiring structural evidence and sufficient interfaces}
\label{sec:acquisition}

Exact emergence theorems can start from complete response laws or incidence data. Acquisition asks how much of the structural presentation can be obtained from feasible observations, and which further tests can resolve competing structures. The distinction between evidence and ideal laws is central from the start. Statistical acquisition guarantees can be staged without requiring every exact theorem to be a practical learner; learning dynamics is developed as a prominent candidate in Section~\ref{sec:learning}.

\subsection{Separate the access models}

\begin{enumerate}
\item \textbf{Supplied exact presentation.} A finite generator, response oracle, or complete algebraic presentation is given. Exact symbolic conclusions are relative to that input. A finite table without a model-class promise is not a complete infinite-horizon specification.
\item \textbf{Exact query learning.} Membership, equivalence, counterexample, reset, or intervention oracles are available only when explicitly granted. Their computational and physical costs must be stated. Existing categorical query-learning theorems are the first reduction target, not new results to rename \cite{CALF,FunctorialLearning}.
\item \textbf{Finite sampling.} Protocols produce random data, potentially with adaptive choice and dependence. Specify resetability, mixing or trajectory assumptions, accessible actions, horizon, observation noise, and a confidence level. Confidence sets must be justified under that sampling law.
\item \textbf{Resource-bounded distinguishability.} Only probes within a time, memory, horizon, or sample budget are accessible. Distinct mathematical behaviors may then remain indistinguishable. This can be a legitimate target of a compositional abstraction theory, including a later computational or cryptographic instance, rather than an inconvenient approximation to unrestricted equivalence.
\end{enumerate}

\subsection{The acquisition target}

Given observational data and a stated class or inference promise, seek an adaptive procedure which produces candidate objects, relations, operations, or interfaces from Target~\ref{target:central}, with certificates appropriate to its identification claim. When alternatives remain, it should identify an executable discriminating experiment or report an unresolved identified set with a justified lower bound. For a specified downstream use, a sufficiency certificate is one possible output. Finite termination or a finite sample bound needs assumptions such as finite distinguishability, a bounded model class, excitation/observability, a positive separation margin, or stability. They must be stated and, where possible, independently tested. A finite rank upper bound learned only by fitting a finite table is not a verified promise about all future experiments.

The procedure should be driven by an unresolved distinction, relation, factorization, or failed certificate in the exact construction: a missing algebraic relation test, an unresolved continuation, a separating positive effect, a joint-correlation probe, or a specific intervention. This is where categorical structure could do work beyond packaging, by transporting a known adequate experiment across modules or localizing the reason an observational interface is insufficient. A proof must quantify the advantage or limitation actually obtained; no blanket efficiency advantage over existing approaches is assumed.

Exact native minimum size is generally unstable under perturbations. A sampled estimate should not be rounded into an exact rank or zero/nonzero support claim without a margin. A more defensible output is a confidence family of resource profiles,
\[
 \left\{\mathscr B_{d,C}(H):H\in\calU_E\right\},
\]
with lower and upper conclusions whose quantifiers are explicit. A theorem about a particular optimum requires additional separation from the relevant realization boundary. A posterior or regularizer can choose a model inside the remaining ambiguity; it does not convert that choice into identification.

\subsection{Approximate composition}

For a chosen operational discrepancy, seek an error law for the declared wiring and interventions. A schematic target is
\[
 d_W\bigl(\Obs(WX),\Obs(W\widehat X)\bigr)
 \leq \omega_W(\varepsilon_1,\ldots,\varepsilon_m;\text{declared stability data}),
\]
where each local approximation has its own certified error. The modulus must follow from actual system and wiring properties. For feedback, conditioning on rare events, or long horizons it may grow badly or fail uniformly. A generic triangle inequality is a baseline; the useful theorem must identify the structural hypotheses making the bound informative and construct examples showing their necessity.

A no-go lower bound should use two admissible models that are indistinguishable under the allowed probes but differ on the structural use. Passive agreement with different intervention responses, finite-horizon agreement with later divergence, and locally indistinguishable components with different joint behavior are mandatory controls. A later experiment is informative only if it is actually executable under the declared access model.

\section{Resource laws and computational methods}
\label{sec:resources}

Resource theory belongs in the programme because a structure can be adequately specified but infeasible to obtain or realize. The desired laws should distinguish:
\begin{enumerate}
\item the cost of the canonical abstract or predictive representation;
\item the minimum cost among native implementations of the observable behavior;
\item the extra cost, or impossibility, of retaining chosen mechanisms or intermediate witnesses;
\item the information and execution cost of using the interface in a composite;
\item the experiment and computation required to certify any of these claims.
\end{enumerate}
The finite hidden/infinite predictive example and the unbounded joint-correlation-memory family show that these distinctions are substantive.

Under a lawful witness map $t$, the inherited value
\[
 (t_!c)(z')=\inf\{c(z):t(z)=z'\}
\]
has a standard associative elimination law. That identity is not a new central theorem and does not make $t_!c$ equal to an independently defined target cost \cite{Interface,Variational,SteinSamuelson}. A native-sufficiency certificate can prove equality, a controlled gap, or impossibility for the latter where appropriate. It should also say when coarse freedom creates new witnesses rather than merely permits refactoring a transported one.

Optimization and compositional gradient methods are possible computational tools. The audited hypergraph result gives a correct assembly of local gradient contributions through linear variable-sharing maps; it does not guarantee existence, global convergence, native interpretation of off-zero iterates, or forward trajectory transport under observation erasure \cite{Budget,Interface,Hanks}. The exact quartic and its counterexamples remain regression tests for any proposed optimizer bridge. A new parameterization, metric, or constraint-preserving dynamics is a separate choice requiring its own equivalence and convergence claims.

The computational obligation is modest in one respect and strict in another. General polynomial-time synthesis is not presumed. Exact witness verification, exact lower-bound or infeasibility certificates when claimed, bit complexity, and clearly reported failure of numerical methods are required. A practical heuristic can be useful while an exact problem remains open, but it must not be substituted for a theorem about native feasibility.

\section{Generality through unlike mathematical instances}
\label{sec:generality}

The theory should transcend a motivating example in the way the contextad work did. Applications provide independent tests of the same emergence and comparison mechanism, not separate definitions of phenomenology or a promise to solve all applications at once.

\begin{enumerate}
\item \textbf{Algebraic and logical observations.} Boolean response profiles, FCA or another generated algebraic doctrine. The test is whether the theorem distinguishes observed-row separation from preservation of generated relations and warranted operation use. A merely stochastic reinterpretation of this example is insufficient.
\item \textbf{Controlled stochastic processes.} Predictive behavior, hidden native implementations, chosen statewise comparisons, sequential wiring, and exposed interventions. The multistage milestone is one demanding native test, including finite hidden models with infinite predictive representations.
\item \textbf{Linear and positive systems.} Finite-dimensional behavioral reconstruction, invariant positive cones, latent realizability, and active identification under declared control and stability assumptions. The theorem must distinguish linear sufficiency from positive and resource sufficiency. Classical linear system identification provides a substantial positive baseline \cite{ActiveID}.
\item \textbf{A materially different later application.} Select one after the central emergence mechanism is clearer: experimental physics with restricted instruments; transactional or robotic systems with discovered operational semantics and partial undo; scientific model composition; or computationally bounded observational equivalence. The selected case must contribute an independent failure mode or theorem, not decorative terminology.
\end{enumerate}

A quantum or physical example must use its own admissible instruments, positivity, and operational access, not ordinary probability with renamed states. A transactional example must specify which changes are available and what partial undo preserves. A resource-bounded computational example must state the adversary/probe class and the intended compositional guarantee. None licenses physical ontology or causal claims from an abstract universal property alone.

Before claiming generality, instantiate the same substantive theorem in at least the algebraic and stochastic/positive settings, and work through a finite-dimensional linear setting or explain precisely which additional analytic assumptions it requires. Examples may serve different roles: a small counterexample isolates a boundary, whereas a parametrized family or real benchmark tests whether inference difficulty remains. The proposal requires both roles and keeps them distinct.

\section{Mandatory stress tests}
\label{sec:stress}

For each proposed criterion or comparison, give an applicability ledger for the following suite: test every control relevant to the claimed semantics and state why the others lie outside its scope. A pure incidence theorem need not settle feedback, and an exact-law theorem need not provide a sampling guarantee. Claims of whole-programme generality must eventually cover the full suite. Existing answers are controls; reproducing them is not a new milestone.

\begin{enumerate}
\item \textbf{A structure must not be hidden in the input.} Re-run the supplied/derived ledger after changing the inference doctrine or hypothesis grammar. State which structural output changes and why. A pairing can be fixed while meet completion, linear span, and native factorization ask different questions.
\item \textbf{Evidence is not access or a full law.} More samples from old tests, newly admissible tests, and a hypothetical enlarged law must produce separately typed revisions. A confidence family must not be called an exact law or assumed nested without justification.
\item \textbf{A learning representation can conceal future distinctions.} The scalar $w=uv$ control has equal present predictions and unequal training responses. A claimed sufficient learning state must close under the declared protocols, while a stationary loss and restricted probes test the scope of that claim.
\item \textbf{Coordinate changes are not new structure.} Transport both dynamics and response maps when testing a gauge change. Reinitializing or imposing a new Euclidean gradient in new coordinates is a different experiment, not a pure invariance check.
\item \textbf{A learned decomposition can be preinstalled.} Compare the proposed output with architecture, regularizer, vocabulary, and initialization controls. State whether the result is identification, selection, or fitting of a designated form.
\item \textbf{Non-equivalence is not universal non-comparability.} A failed inverse or adjoint must not erase a valid one-sided comparison. The allowed comparison class and preserved semantics must be fixed before a no-go claim.
\item \textbf{Observation alone need not identify operations.} Even separating profiles on a finite set do not specify a binary multiplication when no composite observations or algebraic doctrine is supplied. Such a negative control locates the input boundary; it is not the proposed new theorem.
\item \textbf{Emergence can be noncanonical.} Different latent decompositions can have the same observation law. A prior or compression score can select one without identifying it. Seek a structural invariant, an equivalence class, or a specific additional probe rather than silently treating one choice as forced.
\item \textbf{Completion creates new distinctions.} Boolean rows $10$ and $01$ remain distinct after retaining the first coordinate, while their meet/top completion loses distinctions. The criterion must not equate adequacy for observations with adequacy for generated structure.
\item \textbf{Native geometry is not the abstract hull.} Affinely isomorphic square response images can have different minimum latent costs in different output embeddings. The theorem must retain the admissible response embedding and latent-profile class.
\item \textbf{Predictive size and hidden size differ.} A two-state hidden bit model has infinitely many predictive states. Hiding output can destroy the map on previous predictive classes. A new filter or a set-valued comparison may be necessary.
\item \textbf{Joint preparation information is missing locally.} Individually rank-one behavior can have joint rank two under correlated preparation. The theorem must not reconstruct that joint behavior from local reductions alone.
\item \textbf{A chosen mechanism can block refinement.} Early/late randomness has equal coarse trace laws but different compatibility with a newly exposed signal under a specified statewise comparison. More hidden states do not automatically repair the anchor.
\item \textbf{Full support does not imply quantitative liftability.} The two-stage full-support controls share a Boolean possibility pattern but differ in stochastic feasibility and native minimum size. Support-level adequacy is not numerical sufficiency.
\item \textbf{One-step joint feasibility does not settle sequential assembly.} Independently chosen intermediate completions may have no simultaneous cover. A fixed global tensor makes rebracketing coherent but does not solve this stronger problem.
\item \textbf{Local memory does not bound joint memory.} Fixed one-state local factors can require unbounded joint correlation memory. All correlation carriers, phase clocks, and shared random resources must remain visible.
\item \textbf{Value equality does not select an optimum.} Symmetry can prevent an equivariant minimum-state choice. A set of optimal objects or a resource correspondence may be the correct output.
\item \textbf{Feasibility transport does not transport the solver.} Actual quartic zero sets, residual norms, gradients, and global minimum values behave differently under interface changes. A new solver theorem must confront the exact counterexamples.
\item \textbf{Finite information leaves alternatives.} Models can agree on all tested protocols and disagree after an untested continuation or intervention. The claimed acquisition result must use its class promise and access model here.
\item \textbf{Feedback can amplify error.} An approximate compositional claim must display stability or other control, or report a rigorous loss of a uniform bound.
\end{enumerate}

A criterion need not make every test positive. It must return the right kind of negative or noninvertible answer without concealing the loss by changing the native problem. If all difficult tests are excluded by its assumptions, the claimed generality must be reduced accordingly.

\section{Stages and decision gates}
\label{sec:gates}

These are proposed acceptance criteria, not accepted obligations or a schedule. The first decision is to select a concrete theorem on a tractable but nontrivial class, stating its exact positive conclusion, source-level residual, counterexample tests and honest stop condition. It should serve the comparative goal, the generative goal, or both; any unserved programme-level commitment remains open. This draft has not made that selection and is not yet a theorem-ready contract. A single theorem can meet both goals if each conclusion is explicit. The entire native and empirical programme need not be completed before an intermediate result is credited.

The dependency structure is
\[
 \text{input and source audit}
 \longrightarrow
 \begin{cases}
 \text{comparative positive bridge},\\
 \text{generative positive construction},
 \end{cases}
 \longrightarrow
 \text{native, empirical, and cross-instance tests}.
\]
Learning dynamics is a prominent candidate for the generative branch and must be evaluated as such, rather than imposed as its mandatory benchmark. Acquiring finite-sample guarantees is a further task with additional assumptions. A failure in either central branch cannot be compensated for by producing more downstream manuscripts.

\begin{description}[style=nextline,leftmargin=0pt]
\item[Gate A \quad State the first research problem]
Produce the supplied/derived ledger for the selected instances and a source comparison with incidence/FCA, predictive reconstruction, CALF, realization, Stacey, and the existing cylinder work. State the actual theorem to be attempted, its semantics, class of inputs, positive conclusion, independent assumptions, and one concrete premise/conclusion obligation not already obtained from the strongest predecessor. Fix what makes the selected mechanisms unlike, whether the boundary is an if-and-only-if criterion, a hypothesis-removal counterexample or matching quantitative bounds, and whether the claimed tractability is effectiveness, decidability, a complexity bound or experimental feasibility. Compare candidate organizing languages rather than assuming one wins. \emph{Pass}: a tractable theorem-sized question with a genuine residual and explicit controls. \emph{Fail}: a generic fibre or wrapper, a known theorem renamed, no stated experimental interpretation, or hypotheses containing the answer. Reformulate before enlarging the framework.

\item[Gate B \quad Establish a generative result]
Meet Target~\ref{target:generative}, including its positive construction and all six nonvacuity requirements. A response-defined learning-state theorem is a prominent candidate (Target~\ref{target:learning}); an unsupplied interface or decomposition is the demanding flagship (Section~\ref{sec:flagship}). Neither is declared proved or compulsory as the sole form of output. \emph{Pass}: a new theorem derives or explains selection of a useful structural invariant, presentation, or identified family on a nontrivial class, with an exact boundary and a substantive relation to another reconstruction mechanism. \emph{Fail}: all claimed output is designated in advance, parameters are merely renamed, training loss alone is the evidence, the construction only repeats known image/quotient machinery, or no positive family survives. A negative result is recorded as such; it does not automatically pass this gate.

\item[Gate C \quad Establish a comparative result]
Meet a precisely stated first case of Target~\ref{target:central}, with a positive cross-mechanism bridge and an explicit loss or obstruction boundary. The full comparative claim is promoted only after incidence/closure, predictive, and native-factorization mechanisms have been tested at their correct types. \emph{Pass}: independently interpretable changes of data, tests, doctrine, or learning regime explain agreement or a directed comparison and its limitations beyond known preservation by uniqueness. \emph{Fail}: only a shared diagram, arbitrary common embedding, or generic factorization theorem connects the examples. A non-comparison theorem must specify its quantified class and semantics; one failed functor is insufficient.

\item[Gate D \quad Establish native use where claimed]
For the interface flagship, retain Targets~\ref{target:multistage} and~\ref{target:finite}, including stochastic links, sequential reuse, independently constructed assembly routes, intervention semantics, and resource bounds. Relate the result to the existing cylinder comparisons. \emph{Pass}: a new positive native theorem with a sharp obstruction/repair boundary. \emph{Fail}: only rebracketing a supplied full witness, static marginal extension, or a known semipullback is reissued as the new result. A generative theorem of another kind can pass Gate B while this flagship remains explicitly open; it cannot claim the flagship's conclusions.

\item[Gate E \quad Establish the claimed empirical warrant]
Name the exact-query or sampling regime needed for the claimed discovery result, rather than merging them. Link an acquisition, identifiability, or lower-bound theorem to the generated structure. \emph{Pass}: executable probes, sampling/access assumptions, uncertainty or termination guarantees, ambiguity, and resource costs are explicit and nonvacuous. \emph{Fail}: a full-law oracle is called finite evidence, a hypothetical probe is called admissible, prior selection is called identification, or exact boundaries are inferred without stability assumptions. Exact mathematics can be credited before this gate, but claims of feasible empirical discovery cannot.

\item[Gate F \quad Test transfer and assess the thesis]
Apply the same substantive comparison or generation law to unlike instances, including static and sequential data, and, if learning dynamics is selected, assess its actual reconstruction or formation claim rather than a renamed gradient assembly. \emph{Pass}: the common mathematics proves a consequence not supplied by the separate baselines; restrictions and failed transfers remain explicit. \emph{Fail}: only notation transfers. State whether the supported scope is one mechanism, a related family, or a restricted classification. Do not claim a universal theory on this evidence alone.
\end{description}

\subsection{Stop conditions and changes of direction}

For each attempted theorem, fix the input class, positive conclusion, evidence/access regime, and required comparison before evaluating it. Stop claiming that theorem when a verified counterexample defeats its conclusion under the stated hypotheses, the source audit reduces the whole claim to known work, or every surviving condition presupposes the sought solution. Numerical failure is not a proof of impossibility. An inaccessible required probe is a blocker to the empirical claim, not a reason to silently assume an oracle.

Then record one of three outcomes: the target is proved under its original conditions; it remains open with the exact gap identified; or it is refuted in its stated form. A proposed repair must state what changed and preserve a useful positive claim. A restricted theorem is credited at its restricted scope. A negative classification may motivate a new research contract, but cannot retroactively fulfil a positive commitment without an explicit decision to change that commitment.

A successful revision of the contract is not adoption of any of these research stages. The present deliverable stops at revising and critically assessing the proposal. No source audit, construction, experiment, or ongoing research beyond that task is promised by these proposed gates.

\subsection{Current gaps and dependencies}

\begin{enumerate}
\item \textbf{Organizing language.} Two-sided testing is a supported candidate. No theorem yet establishes its necessity, sufficiency for the full scope, or superiority to every alternative. Dependence: faithful instance presentations and a nontrivial comparative result.
\item \textbf{New comparative bridge.} The existing row/completion, generated-image, cylinder, positive, and native results remain baselines. Dependence: one precisely stated residual claim beyond them; all further comparisons remain open until proved.
\item \textbf{Generative theorem.} Broad targets are proposed, not constructions already obtained. Dependence: a specified output and nonvacuous input-level criterion. The elementary learning-flow example motivates the issue but does not discharge it.
\item \textbf{Unsupplied decomposition and multistage use.} The earlier positive targets are retained and open. Dependence: native admissibility, joint correlations, specified stochastic comparisons, and legal sequential/intervention semantics.
\item \textbf{Evidence to structure.} Exact-law results do not by themselves identify structures from records. Dependence: justified access and sampling models, identifiability or an identified set, and quantitative guarantees when claimed.
\item \textbf{Generality and novelty.} Neither a list of applications nor independent review constitutes novelty clearance. Dependence: the same added theorem in unlike cases and a focused strongest-source audit of that theorem.
\end{enumerate}

\section{Proof discipline and deliverables}
\label{sec:discipline}

Each substantive result should be accompanied by:
\begin{enumerate}
\item a full statement of native objects, observations, operations, comparison maps, and all finiteness, positivity, reachability, separation, access, and resource assumptions;
\item a proof in which the independent criterion is separated from the existence conclusion, and a typed diagram showing what commutes or only compares;
\item a strongest-source audit identifying the exact predecessor theorem, the reduction to it, and the additional claim, with full text checked where a precise theorem attribution is made;
\item exact examples, obstruction certificates, and reproducible finite checks for the claimed formulas, including zero masses, unused states, empty fibres, and infinite-cost cases;
\item an independent mathematical review and a separate source/provenance review; neither finite testing nor independent review is called formal verification;
\item a status ledger distinguishing proved results, conditional conjectures, counterexamples, repaired theorems, computational evidence, and unresolved tasks.
\end{enumerate}

The recent reports provide a good control standard. Their exact checks and independent reviews support the stated finite formulas; their source audits explicitly avoid novelty clearance. The fixed-budget audit distinguishes coercivity and exact certificates from optimizer success. The interface audit checks the actual residual family and variance, rather than substituting a convenient different objective. These standards should carry forward.

Suggested outputs are an emergence/comparison theorem manuscript, a chapter relating native realization and refinement to that construction, a source-and-boundary ledger, and a small exact verification suite. An acquisition paper follows only when its access and statistical assumptions are adequate. The deliverables are mathematical results, not a count of manuscripts or definitions. No deadline, publication commitment, implementation commitment, or adoption of this proposal is implied by this draft.

This deliverable is intentionally uncompiled \LaTeX{} source. Static checks cover source structure and cross-references; no PDF layout or compilation is claimed. The proposed new theorems remain research targets, and this draft does not adopt them as a contract.

\section{What the proposal would let us conclude}
\label{sec:conclusion}

If successful, the programme would explain how a structure arises from observational patterns before it is treated as a representation to preserve. It would relate generated distinctions, concepts, predictive states, induced operations, and latent or decompositional explanations, while distinguishing their warrants and choices. Its central contributions would explain a substantive case of structure generation or learning-driven selection, and why different reconstruction mechanisms agree, compare with loss, or resist a specified comparison. Candidate interfaces remain one important output rather than a universal requirement. Native realization, refinement, and resources would establish what those structures can support.

If the strongest unification fails, a classification of the unavoidable divisions could still establish a fundamental architecture for these problems: what observations force, what realization adds, where choice enters, and which changes have only one-sided or resource-sensitive comparisons. The failure must be proved at the appropriate level, not inferred from one inconvenient example.

The proposed next decision is whether the revised comparative and generative targets, their independent acceptance tests, and the retained flagship capture the desired research direction. Only then should a particular next theorem or research stage be selected. The supplied results already justify taking the question seriously. They do not yet establish that categorical phenomenology is one canonical construction, a new mathematical discipline, or a universal method of discovering primitives. The proposal is designed to find out which stronger claim can actually be defended.

\appendix
\section{Assumptions that must remain visible}
\label{app:assumptions}

\begin{enumerate}
\item A repertoire of measurements/actions is supplied at each stage. Revision is allowed, but a new primitive needs an interpretation or an explicitly conjectural status.
\item Algebraic closure and native admissibility are separate inputs. Formal joins, convex mixtures, tensor products, or compact completions need not be executable states or operations.
\item A named native comparison has a direction and strength. Public traces, conditional kernels, statewise simulation, deterministic projection, and intervention preservation are not interchangeable.
\item A finite table is not complete behavior without a finite-presentation or model-class assumption. A finite-dimensional fit is not an independently established dimension bound.
\item Declared hidden commitments are not automatically empirically identified. A specified statewise kernel can be retained as a design commitment while remaining observationally underdetermined.
\item Local observational adequacy does not give joint correlation adequacy. Allowed correlations, shared randomness, preparation scopes, and wire interfaces are explicit.
\item Unbounded native existence, fixed-budget feasibility, exact optimum, approximate fit, and optimizer convergence are separate questions.
\item Budget preservation for every coarse model is stronger than existence, equality of optimum values, or comparison of chosen endpoints.
\item Noisy support, rank, and exact resource boundaries require margins or set-valued conclusions. Conditioning and feedback require their own error control.
\item Every enlargement of a model class or weakening of commitments changes the problem unless a comparison theorem proves the intended original use survives.
\end{enumerate}

\section{Source register and scope of attribution}
\label{app:sources}

The six local reports and their review files are the primary record of the recent work. The recovered ledger identifies the earlier Boolean/stochastic, predictive, positive-Hankel, cylinder, and cross-doctrine checkpoints. The draft uses those sources to set the baseline; it does not claim to have independently reproved every recovered exploratory theorem.

The external references below are deliberately selective. The emergence precedents in Section~\ref{sec:central} control claims of structural, feature, and operation discovery; Garner controls the generic cylinder/Isbell claims; functorial automata and categorical learning control generic reconstruction and exact acquisition; statistical comparison and causal abstraction control the experimental and intervention language; positive realization and multicausal coupling control native and temporal claims; variational analysis and the gradient paper control optimization. Their juxtaposition is not evidence that the proposed bridge is absent from the literature. Before a novelty assertion, the precise target must be checked against the full strongest sources, including their hypotheses and later developments.

McCullagh's design-indexed statistical models, Lorenz--Tull's query-indexed causal abstractions, and double-operadic systems doctrines are especially important controls: changing experimental design, naturality of abstraction, and compositional interface syntax already have substantive theories. The proposed emergence theorem must coexist with these theories, and its refinement component must add a specific structural or native consequence rather than merely reproduce their languages \cite{McCullagh,LorenzTull,SystemsDoctrines}.

\begin{thebibliography}{99}
\bibitem{TwoSidedNote}
\emph{Towards two-sided context-sensitive factorization-based realization}.
Inspected programme note. Used for its actual constructions, comparisons, and stated limits; no original authorship date is inferred from the file's current metadata.

\bibitem{Stacey}
Andrew Stacey. \emph{Comparative Smootheology}.
\url{https://arxiv.org/abs/0802.2225}.
Numbering in Section~\ref{sec:stacey} refers to the inspected arXiv PDF, not to a secondary summary.

\bibitem{Core}
\emph{A common generated observable construction and its native boundaries}.
Programme report, 2 November 2026. Accompanied by source audit, recovered theorem ledger, exact checks, and review.

\bibitem{Refinement}
\emph{A focused reduction test for refinement of a chosen stochastic realization}. Programme report, 2 November 2026.
Finite controlled stochastic refinement, exact real-algebra reduction, two-stage garbling and partition criteria. Accompanied by source audit and review.

\bibitem{Common}
\emph{A common native construction for finite relational and stochastic refinements}.
Programme report, 2 November 2026. Accompanied by independent lifting criteria, source audit, stress tests, and review.

\bibitem{Joint}
\emph{Joint native realization with fixed local mechanisms}.
Programme report, 3 November 2026. Finite tensor and row-polytope assembly, constrained-rank resources, witness-fibre initiality, and transport. Accompanied by source audit and review.

\bibitem{Budget}
\emph{Fixed budget native realization and compositional gradient assembly}.
Programme report, 3 November 2026. Accompanied by exact residual checks, source audit, and mathematical review.

\bibitem{Interface}
\emph{Interface changes and fixed budget optimization}.
Programme report, 3 November 2026. Accompanied by exact examples, source audit, and independent review.

\bibitem{Recovered}
Recovered theorem ledger, 2 November 2026.
Records the earlier Boolean and stochastic discussions, predictive construction, positive Hankel construction, completed cylinder/wiring checkpoint, and completed cross-doctrine checkpoint. Bibliographic and proof status varies by source.

\bibitem{ComputationalMechanics}
Cosma Rohilla Shalizi and James P. Crutchfield.
\emph{Computational Mechanics: Pattern and Prediction, Structure and Simplicity}.
Journal of Statistical Physics 104, 817--879, 2001.
\url{https://arxiv.org/abs/cond-mat/9907176}.

\bibitem{FCA}
Rudolf Wille. \emph{Restructuring Lattice Theory: An Approach Based on Hierarchies of Concepts}. 1982.
\url{https://doi.org/10.1007/978-94-009-7798-3_15}.

\bibitem{ContextualGrammar}
Alexander Clark, R\'emi Eyraud, and Amaury Habrard.
\emph{Using Contextual Representations to Efficiently Learn Context-Free Languages}.
Journal of Machine Learning Research 11, 2707--2744, 2010.
\url{https://jmlr.org/papers/v11/clark10a.html}.
Its identification-in-the-limit and update/query-work bounds do not give unrestricted finite-sample exact recovery.

\bibitem{MacroVariables}
Krzysztof Chalupka, Frederick Eberhardt, and Pietro Perona.
\emph{Multi-Level Cause-Effect Systems}. AISTATS 2016.
\url{https://proceedings.mlr.press/v51/chalupka16.html}.
Used for explicit variable/decomposition construction with supplied causal and intervention assumptions, not an unchecked minimal-statistical-sufficiency claim.

\bibitem{StructuralForm}
Charles Kemp and Joshua B. Tenenbaum.
\emph{The discovery of structural form}. PNAS 105, 10687--10692, 2008.
\url{https://pmc.ncbi.nlm.nih.gov/articles/PMC2492756/}.

\bibitem{DreamCoder}
Kevin Ellis and collaborators.
\emph{DreamCoder: Bootstrapping Inductive Program Synthesis with Wake-Sleep Library Learning}. PLDI 2021.
\url{https://doi.org/10.1145/3453483.3454080}.

\bibitem{Grokking}
Neel Nanda and collaborators. \emph{Progress measures for grokking via mechanistic interpretability}. 2023.
\url{https://arxiv.org/abs/2301.05217}.
Sections 4--6 provide the mechanism, interventions, progress measures, and stated limitations used here.

\bibitem{Disentanglement}
Francesco Locatello and collaborators.
\emph{Challenging Common Assumptions in the Unsupervised Learning of Disentangled Representations}. ICML 2019.
\url{https://proceedings.mlr.press/v97/locatello19a.html}.
The stated nonidentifiability result is model-relative, not an impossibility theorem for all structural discovery.

\bibitem{AbstractCompletion}
Roberto Giacobazzi, Francesco Ranzato, and Francesca Scozzari.
\emph{Complete Abstract Interpretations made Constructive}. MFCS 1998.
\url{https://www.math.unipd.it/~ranzato/papers/mfcs98.pdf}.
Constructive completeness repair under its continuity assumptions.

\bibitem{StrongPreservation}
Francesco Ranzato and Francesco Tapparo.
\emph{Generalized Strong Preservation by Abstract Interpretation}.
Journal of Logic and Computation, 2007.
\url{https://arxiv.org/abs/cs/0401016}.
Forward and backward completeness have different equations and must not be conflated.

\bibitem{RefinementMappings}
Mart\'in Abadi and Leslie Lamport.
\emph{The Existence of Refinement Mappings}.
\url{https://lamport.azurewebsites.net/pubs/abadi-existence.pdf}.
Its existence theorem uses stated machine-closure, continuity, and nondeterminism assumptions and can add history and prophecy variables.

\bibitem{InterfaceAutomata}
Luca de Alfaro and Thomas A. Henzinger.
\emph{Interface Automata}. 2001.
\url{https://web.cs.wpi.edu/~heineman/html/teaching_/CS562/p109-de_alfaro.pdf}.
Compositional substitution requires its explicit compatibility conditions.

\bibitem{Garner}
Richard Garner. \emph{The Isbell monad}.
\url{https://arxiv.org/abs/1410.7108}.
Cylinder-factorization pseudoalgebras and lax/colax comparisons; not a minimum-native-realization theorem.

\bibitem{Automata}
Thomas Colcombet and Daniela Petri\c{s}an.
\emph{Automata Minimization: a Functorial Approach}.
\url{https://arxiv.org/abs/1711.03063}.

\bibitem{CALF}
Gerco van Heerdt, Matteo Sammartino, and Alexandra Silva.
\emph{CALF: Categorical Automata Learning Framework}.
\url{https://arxiv.org/abs/1704.05676}.
See also the algebraic/tree-automata development at
\url{https://arxiv.org/abs/2001.05786}.

\bibitem{FunctorialLearning}
Thomas Colcombet, Daniela Petri\c{s}an, and Riccardo Stabile.
\emph{Learning automata and transducers: a categorical approach}.
\url{https://arxiv.org/abs/2010.13675}.
Exact-query assumptions and termination hypotheses must be retained.

\bibitem{BlackwellCategory}
Tobias Fritz, Tom\'a\v{s} Gonda, Paolo Perrone, and Eigil Fjeldgren Rischel.
\emph{Representable Markov Categories and Comparison of Statistical Experiments in Categorical Probability}.
Theoretical Computer Science 961, 113896, 2023.
\url{https://arxiv.org/abs/2010.07416}.

\bibitem{Positive}
Mike Boyle and Karl Petersen.
\emph{Hidden Markov Processes in the Context of Symbolic Dynamics}.
Author manuscript, including the positive-realization account and Heller criterion.
\url{https://math.umd.edu/~mboyle/papers/BoylePetersen07-10-2009.pdf}.

\bibitem{Multicausal}
Beatrice Acciaio, Daniel Kr\v{s}ek, and Gudmund Pammer.
\emph{Multicausal transport: barycenters and dynamic matching}.
SIAM Journal on Financial Mathematics 16(3), 1104--1138, 2025.
\url{https://arxiv.org/abs/2401.12748}.
Successive kernel and gluing results are prior art; latent-cardinality claims require their own proof.

\bibitem{McCullagh}
Peter McCullagh. \emph{What is a statistical model?} 2002.
\url{https://www.stat.uchicago.edu/~pmcc/pubs/AOS023.pdf}.

\bibitem{LorenzTull}
Robin Lorenz and Sean Tull. \emph{Causal and Compositional Abstraction}, 2026.
\url{https://arxiv.org/abs/2602.16612}.
Query-indexed and mechanism-level causal abstraction, including interventions; precise attribution is subject to the target's source audit.

\bibitem{SystemsDoctrines}
Sophie Libkind and David Jaz Myers. \emph{Towards a double operadic theory of systems}, 2025.
\url{https://arxiv.org/abs/2505.18329}.
Used as a compositional-syntax control, not as an empirical identification theorem.

\bibitem{ActiveID}
Andrew Wagenmaker and Kevin Jamieson.
\emph{Active Learning for Identification of Linear Dynamical Systems}.
COLT 2020.
\url{https://proceedings.mlr.press/v125/wagenmaker20a.html}.

\bibitem{Variational}
R. Tyrrell Rockafellar and Roger J.-B. Wets.
\emph{Variational Analysis}. Springer, 1998.
\url{https://sites.math.washington.edu/~rtr/papers/rtr169-VarAnalysis-RockWets.pdf}.
Equation 1(17) is the pertinent value-function/infimization construction.

\bibitem{SteinSamuelson}
Dario Stein and Richard Samuelson.
\emph{Towards a Compositional Framework for Convex Analysis (with Applications to Probability Theory)}.
FoSSaCS 2024.
\url{https://arxiv.org/abs/2312.02291}.
The native rank and quartic objectives need not satisfy its convexity hypotheses.

\bibitem{Hanks}
Tyler Hanks, James Fairbanks, and Matthew Klawonn.
\emph{Generalized Gradient Descent is a Hypergraph Functor}.
EPTCS 429, 217--233, 2025.
\url{https://arxiv.org/abs/2403.19845}.
Theorems 3.4--3.5 and Corollary 3.5.1 concern field assembly and gradient-flow functoriality, with their additional hypotheses.
\end{thebibliography}
\end{document}
